\section*{Chapter \ref{ch:g9:logic-gates} --- Logic Gates and Boolean Algebra}

\begin{solution}{exo:g9:logic-gates:1}
For $a = 1$, $b = 0$: the OR gives 1, the NOT turns it into 0. For $a = 0$,
$b = 0$: the OR gives 0, and the \omterm{def:g1:what-is-a-computer:input}{output} is 1. The circuit \omterm{def:g1:what-is-a-computer:input}{outputs} 1 only when
neither \omterm{def:g1:what-is-a-computer:input}{input} is 1.
\end{solution}

\begin{solution}{exo:g9:logic-gates:2}
\code{0 0 | 0}, \code{0 1 | 1}, \code{1 0 | 1}, \code{1 1 | 0}: the \omterm{def:g1:what-is-a-computer:input}{output} is
1 when the two \omterm{def:g1:what-is-a-computer:input}{inputs} are different.
\end{solution}

\begin{solution}{exo:g9:logic-gates:3}
\omoutput{code/grade-9/02-logic-gates/outputs/nand_only.txt}{NAND, and the three
gates built from it.}\label{out:g9:logic-gates:nand}

The AND and OR columns are those of the real gates, and NOT $a$ is the opposite
of $a$.
\end{solution}

\begin{solution}{exo:g9:logic-gates:4}
With \omterm{def:g5:bits-and-bytes:bit}{bits}, ``exactly one is 1'' means ``their sum is 1'':

\omcode{code/grade-9/02-logic-gates/exercises/exactly_one.py}{1}{3}{Exactly one of
three.}\label{lst:g9:logic-gates:sol-one}
\end{solution}

\begin{solution}{exo:g9:logic-gates:5}
\code{range(1)} holds only 0, so both \omterm{def:g4:algorithms-repeat-and-decide:loop}{loops} make a single turn and print
\code{0 0 0}. With \code{range(2)}, or \code{[0, 1]}, they print the four
rows.
\end{solution}

\begin{solution}{exo:g9:logic-gates:6}
The \omterm{def:g1:what-is-a-computer:input}{input} $a = 1$, $b = 1$: OR gives a sum \omterm{def:g5:bits-and-bytes:bit}{bit} 1, and the \omterm{def:g8:binary-and-hexadecimal:carry}{carry} is 1, so the
circuit claims that $1 + 1$ is \code{11} in \omterm{def:g5:bits-and-bytes:binary}{binary}, which is 3.
\end{solution}

\begin{solution}{exo:g9:logic-gates:7}
Saved next to \code{gates.py} (\cref{lst:g9:logic-gates:gates}), the \omterm{def:g4:programs-and-operating-system:program}{program} imports
the three gates it needs:

\omcode{code/grade-9/02-logic-gates/exercises/table_or_not.py}{1}{6}{The table of
\code{(a OR b) AND (NOT c)}.}\label{lst:g9:logic-gates:sol-table}

The \omterm{def:g1:what-is-a-computer:input}{output} is 1 on three rows only: \code{0 1 0}, \code{1 0 0} and
\code{1 1 0}; at least one of $a$ and $b$ is 1, and $c$ is 0.
\end{solution}

\begin{solution}{exo:g9:logic-gates:8}
\omterm{def:g7:functions:function}{Functions} can be passed as arguments, like the sort key of
\cref{ch:g8:project-data-explorer}:

\omcode{code/grade-9/02-logic-gates/exercises/agree.py}{1}{7}{Two circuits compared
on every \omterm{def:g1:what-is-a-computer:input}{input}.}\label{lst:g9:logic-gates:sol-agree}
\end{solution}

\begin{solution}{exo:g9:logic-gates:9}
$2^5 = 32$ rows, and $2^{10} = 1{,}024$ with 10 \omterm{def:g1:what-is-a-computer:input}{inputs}.
\end{solution}

\begin{solution}{exo:g9:logic-gates:10}
\code{0 0}: $a$ AND $b$ is 0, the OR is 0. \code{0 1}: 0, then 0.
\code{1 0}: 0, then 1. \code{1 1}: 1, then 1. The last column, 0, 0, 1, 1, is
the column of $a$. In words: if $a$ is 1, the OR is 1 anyway; if $a$ is 0,
then $a$ AND $b$ is 0 too, and the OR is 0.
\end{solution}

\begin{solution}{exo:g9:logic-gates:11}
Each of the four rows has 2 possible \omterm{def:g1:what-is-a-computer:input}{outputs}: $2^4 = 16$ tables.
\Cref{lst:g9:logic-gates:nand} builds NOT, AND and OR from NAND. Now take any
table. For each row whose \omterm{def:g1:what-is-a-computer:input}{output} is 1, an AND of the \omterm{def:g1:what-is-a-computer:input}{inputs}, or of their NOT,
gives 1 on that row only: for the row $a = 1$, $b = 0$, it is $a$ AND
(NOT $b$). The OR of these pieces gives 1 exactly on the rows wanted. If no
row has \omterm{def:g1:what-is-a-computer:input}{output} 1, $a$ AND (NOT $a$) is always 0. So every table is a circuit of
NOT, AND and OR, and therefore of NAND alone. A test of the chapter also
builds all 16 by combining NAND gates.
\end{solution}

\begin{solution}{exo:g9:logic-gates:12}
A \omterm{def:g9:logic-gates:half-adder}{half adder} adds the units; a second one adds the twos; a third adds the
\omterm{def:g8:binary-and-hexadecimal:carry}{carry} of the units to that. The two carries can never both be 1, so an OR
gives the last \omterm{def:g5:bits-and-bytes:bit}{bit}. The \omterm{def:g3:files-and-folders:file}{file} sits next to \code{gates.py}:

\omcode{code/grade-9/02-logic-gates/exercises/two_bit_adder.py}{1}{13}{Adding two
numbers of two \omterm{def:g5:bits-and-bytes:bit}{bits}.}\label{lst:g9:logic-gates:sol-adder}

A test of the chapter checks $4 \times \text{carry} + 2 s_1 + s_0 = x + y$ for
all 16 pairs of numbers from 0 to 3.
\end{solution}

\begin{solution}{pb:g9:logic-gates:1}
\textbf{1.} The rule's column of the \omterm{def:g4:programs-and-operating-system:program}{program}'s \omterm{def:g1:what-is-a-computer:input}{output} below.

\textbf{2.} \code{0 1 1}, \code{1 0 1}, \code{1 1 0} and \code{1 1 1}: the
rows with at least two 1s.

\textbf{3.} Each judge presses or not, 2 choices each: $2^3 = 8$ patterns.

\textbf{4.} If at least two judges press, some pair of judges both press, and
that pair's AND gives 1; if at most one presses, no pair does.

\textbf{5.} If $a$ and $b$ both press, $a$ AND $b$ makes the buzzer sound. If
neither presses, at most one judge pressed, and $a$ OR $b$ is 0, so the \omterm{def:g1:what-is-a-computer:input}{output}
is 0. If they disagree, there is exactly one press already, and $c$ decides:
$c$ AND ($a$ OR $b$) is $c$.

\textbf{6.} Yes:

\omoutput{code/grade-9/02-logic-gates/outputs/judges.txt}{The rule and the two
circuits.}\label{out:g9:logic-gates:judges}

\textbf{7.} The 8 \omterm{def:g1:what-is-a-computer:input}{inputs} are all the \omterm{def:g1:what-is-a-computer:input}{inputs} there are: agreeing on each of them
is agreeing everywhere, not on a sample.

\textbf{8.} \code{pairs}: three ANDs and two ORs, 5 gates. \code{majority}: two
ANDs and two ORs, 4 gates.

\textbf{9.} Three \omterm{def:g9:logic-gates:not-and-or-xor}{AND gates}, one per pair ($a b$, $b c$, $a c$); an OR joins the
first two, and a second OR joins that with the third.

\textbf{10.} $2^5 = 32$ rows.

\textbf{11.} \code{majority} saves a gate; \code{pairs} is easier to check by
reading. Both are proved right by the \omterm{def:g4:programs-and-operating-system:program}{program}.

\textbf{12.} \textbf{Result.} ``At least two of three'' is
($a$ AND $b$) OR ($b$ AND $c$) OR ($a$ AND $c$): five \omterm{def:g1:what-is-a-computer:input}{two-input} gates, and the
\omterm{def:g4:programs-and-operating-system:program}{program} checks it, against the rule, on all 8 \omterm{def:g1:what-is-a-computer:input}{inputs}; a circuit of four gates
does the same. Five judges need a \omterm{def:g9:logic-gates:gate}{truth table} of 32 rows.
\end{solution}
